\chapter{系统需求分析}

\section{应用场景分析}

本文所设计的平台主要面向高校程序设计课程教学、实验训练和日常在线练习等场景。在课程教学场景中，教师需要通过平台发布编程题目、收集学生提交记录并完成平时作业或实验成绩统计；在自主训练场景中，学生需要通过平台进行题目练习、查看判题结果和复盘错误原因；在集中考核或训练高峰场景中，平台还需要在较短时间内处理大量提交任务，保证结果返回及时、环境一致和过程可追踪。

上述场景具有几个共同特征：一是提交行为具有明显的并发性与波动性，容易在作业截止前或训练期间出现集中峰值；二是判题执行对环境隔离和资源控制有较高要求；三是用户不仅关注最终结果，也越来越关注解释性反馈和学习支持。因此，平台需要同时面向“判题正确”“系统稳定”“反馈可用”三个维度展开设计。

\section{用户角色与业务流程分析}

从平台使用对象来看，系统主要涉及学生、教师和管理员三类用户。学生是平台的核心使用者，其主要行为包括注册登录、浏览题目、在线提交代码、查看评测结果和接收 AI 辅助反馈。教师主要负责课程题目管理、测试数据维护、作业布置和提交记录查看。管理员则更多承担系统层面的用户权限管理、资源配置、服务监控和运行维护等职责。

从业务流程角度看，平台的核心流程可以概括为：用户登录后选择题目并提交代码，系统接收提交请求后生成评测任务，调度服务将任务下发给评测执行节点，执行节点在隔离环境中完成编译、运行与测试用例校验，并将评测结果回传至业务服务；若启用 AI 辅助功能，则在客观评测结果生成后，再由 AI 模块基于代码与错误信息输出解释性反馈，最终统一展示给用户。

\section{功能需求分析}

结合上述应用场景和业务流程，平台应至少具备以下核心功能。第一，用户与权限管理功能，包括注册登录、身份认证、角色区分和权限控制。第二，题库管理功能，包括题目创建、题目分类、题面维护、样例管理和测试数据管理。第三，在线提交与自动评测功能，包括代码编辑或上传、任务提交、结果判定和状态展示。第四，提交记录与统计分析功能，包括历史提交查询、通过情况统计和错误类型统计。第五，AI 辅助反馈功能，用于在判题结束后生成有限度的解释与提示。

需要说明的是，复杂的教师管理、学生管理和后台统计面板虽然存在一定工程价值，但并不是本文的研究重点。本文更关注与“在线评测主链路”直接相关的题库、提交、判题和智能辅助模块，而将复杂后台功能视为支撑性或可选增强功能。

\section{非功能需求分析}

在线编程评测平台的非功能需求对系统质量具有决定性作用。首先是正确性与一致性要求。平台必须保证在统一环境和统一测试条件下得到稳定一致的判题结果，这是系统能够用于课程评价与训练分析的前提。其次是安全性要求。由于平台需要运行用户提交的任意程序，必须通过容器隔离、资源限制和执行权限约束等机制降低对宿主环境的威胁。

其次是性能与扩展性要求。平台需要能够处理一定规模的并发提交，并在任务高峰期保持较稳定的响应和调度能力，因此应具备异步处理、水平扩展和资源弹性调度能力。再次是可维护性与可观测性要求。系统应支持模块化部署、日志记录、运行监控和异常追踪，便于后续调试、优化和故障定位。最后是可用性要求。平台交互界面应清晰、反馈信息应易于理解，尤其是 AI 辅助模块的输出应简洁、克制并对学习过程真正有帮助。

\section{本章小结}

本章围绕平台的典型应用场景、用户角色、核心业务流程以及功能与非功能需求展开分析，明确了系统建设必须兼顾自动评测、运行安全、资源弹性和学习支持等多方面目标。这些需求分析结果将直接作为下一章总体设计的依据。
